\section{潮流求解方法}
\label{sec:methods}

\subsection{牛顿--拉夫逊法}

牛顿--拉夫逊法是潮流方程的经典求解方法\cite{tinney1967}，在每轮迭代中求解
\begin{equation}
\bm{J}(\bm{x}^{(k)})\Delta\bm{x}^{(k)}=-\bm{f}(\bm{x}^{(k)}),\qquad
\bm{x}^{(k+1)}=\bm{x}^{(k)}+\Delta\bm{x}^{(k)}.
\end{equation}
当初值接近非奇异解时，该方法通常具有二次收敛特征，但每轮需要重新构造并分解雅可比矩阵。

\subsection{快速解耦潮流法}

快速解耦法利用 $P$--$\theta$ 与 $Q$--$V$ 的主导耦合关系\cite{stott1974}，将迭代近似分解为
\begin{equation}
\bm{B}'\Delta\bm{\theta}=\Delta\bm{P}/\bm{V},\qquad
\bm{B}''\Delta\bm{V}=\Delta\bm{Q}/\bm{V}.
\end{equation}
固定系数矩阵可以预先分解，因此单轮成本低；但高 $R/X$、强耦合或临界工况会削弱其近似条件。

\subsection{全局化方法与连续潮流}

最优乘子法在牛顿方向上搜索合适步长，非线性规划法最小化 $\tfrac12\|\bm{f}(\bm{x})\|_2^2$。它们可以抑制迭代发散，但也可能停在非零残差的最小二乘点。

连续潮流法采用预测--校正策略跟踪 $\bm{f}(\bm{x},\lambda)=\bm{0}$ 的解流形\cite{ajjarapu1992}，并在鼻点附近使用局部或伪弧长等参数化绕开自然参数化的奇异。项目当前支持自然、伪弧长、切线角和绝对 $V/P$ 角参数化。

\subsection{方法之间的定位关系}

NR 与 FDLF 用于求取给定参数下的离散运行点；最优乘子和非线性规划用于改善定点求解过程的全局性；CPF 则用于追踪随负荷参数连续变化的解分支。前三类方法回答“当前运行点能否求出”，CPF 回答“运行点如何随参数演化并在何处达到极限”。因此，算法比较不能只比较耗时，还应同时报告最终残差、步长控制量、无功限值状态和所处解支。
